\chapter{Moser Iteration and Supersolution Estimates}\label{sec:moser}
%% ========================================================================

%% --- Section 7 ---

\section{Moser Iteration}

Throughout this section we fix a dimension $d \in \N$ with $d \geq 1$ and
$d > 2$, and write $E = \R^d$ for the ambient space. All balls are taken
around the origin in $E$, and we abbreviate $\Bz{r} = \Ball{r}{0}$. Functions
$u : E \to \R$ are real valued; $A$ denotes a normalized elliptic coefficient
field on $\Bz{1}$, with ellipticity ratio $\Lambda = \Lambda_A \geq 1$
\leanname{NormalizedEllipticCoeff}. The notion of subsolution is
\leanname{IsSubsolution} and that of supersolution is
\leanname{IsSupersolution}. The starting point is the normalized De Giorgi
$L^2 \to L^\infty$ estimate proved in Section 5.

\subsection{Moser constants}

\begin{definition}[Moser anchor constant]
\label{def:C_MoserAnchor}
\lean{DeGiorgi.C_MoserAnchor}
\leanok
\uses{def:C_gns, lem:smoothTransition_nnnorm_deriv_bounded, def:K_DeGiorgi_subsolution_normalized}
The \emph{Moser anchor} in dimension $d$ is
\begin{equation*}
  C_{\mathrm{MA}}(d) := \max\!\left\{1,\;
    C_{\mathrm{DG}}^{\mathrm{sub}}(d)^{2},\;
    8\,\Mst^{2}\,C_{\mathrm{gns}}(d,2)^{2}\right\},
\end{equation*}
where $C_{\mathrm{DG}}^{\mathrm{sub}}(d)$ is the normalized De Giorgi
subsolution constant of Section 5 and $C_{\mathrm{gns}}(d,2)$ is the
Gagliardo--Nirenberg--Sobolev constant for $p=2$.
\end{definition}

\begin{lemma}
\label{lem:one_le_C_MoserAnchor}
\lean{DeGiorgi.one_le_C_MoserAnchor}
\leanok
\uses{def:C_MoserAnchor}
$1 \leq C_{\mathrm{MA}}(d)$.
\end{lemma}

\begin{proof}
\leanok

Immediate from $C_{\mathrm{MA}}(d) \geq \max\{1,\dots\} \geq 1$.
\end{proof}

\begin{definition}[Sobolev gain factor and decay ratio]
\label{def:moserChi}
\lean{DeGiorgi.moserChi, DeGiorgi.moserDecayRatio}
\leanok
Set
\begin{equation*}
  \chi := \frac{d}{d-2} \quad (\text{\leanname{moserChi}}),
  \qquad
  q := \frac{d-2}{d} \quad (\text{\leanname{moserDecayRatio}}).
\end{equation*}
\end{definition}

\begin{lemma}[Basic algebraic facts on $\chi,q$]
\label{lem:moserDecayRatio_nonneg}
\lean{DeGiorgi.moserDecayRatio_nonneg, DeGiorgi.moserDecayRatio_lt_one, DeGiorgi.moserChi_pos, DeGiorgi.one_lt_moserChi, DeGiorgi.moserDecayRatio_eq_inv_moserChi}
\leanok
\uses{def:moserChi}
Assume $d > 2$. Then
\begin{enumerate}
  \item $0 \leq q$ \leanname{moserDecayRatio\_nonneg};
  \item $q < 1$ \leanname{moserDecayRatio\_lt\_one};
  \item $0 < \chi$ \leanname{moserChi\_pos};
  \item $1 < \chi$ \leanname{one\_lt\_moserChi};
  \item $q = \chi^{-1}$ \leanname{moserDecayRatio\_eq\_inv\_moserChi}.
\end{enumerate}
\end{lemma}

\begin{proof}
\leanok

For (1)--(2) note $0 \leq d-2 < d$, so $0 \leq (d-2)/d < 1$.
For (3)--(4), $d/(d-2) > 0$ and $d > d-2$ gives $d/(d-2) > 1$.
For (5), $\big((d-2)/d\big)^{-1} = d/(d-2)$.
\end{proof}

\begin{definition}[Dimensional Moser constant]
\label{def:C_Moser}
\lean{DeGiorgi.C_Moser}
\leanok
\uses{def:C_MoserAnchor, def:moserChi}
Set $C_{\mathrm{MA}} := C_{\mathrm{MA}}(d)$. If $d > 2$, define
\begin{equation*}
  C_{\mathrm{Moser}}(d) := \max\!\left\{C_{\mathrm{MA}},\;
    (32\,C_{\mathrm{MA}})^{d/2}\,
    4^{\sum_{n=0}^{\infty} n\,q^{n}}\right\};
\end{equation*}
otherwise $C_{\mathrm{Moser}}(d) := C_{\mathrm{MA}}$.
\end{definition}

\begin{lemma}
\label{lem:C_MoserAnchor_le_C_Moser}
\lean{DeGiorgi.C_MoserAnchor_le_C_Moser, DeGiorgi.one_le_C_Moser, DeGiorgi.C_DeGiorgi_subsolution_normalized_sq_le_C_Moser}
\leanok
\uses{def:K_DeGiorgi_subsolution_normalized, def:C_MoserAnchor, def:C_Moser}
$C_{\mathrm{MA}}(d) \leq C_{\mathrm{Moser}}(d)$, $1 \leq C_{\mathrm{Moser}}(d)$, and
$C_{\mathrm{DG}}^{\mathrm{sub}}(d)^2 \leq C_{\mathrm{Moser}}(d)$.
\end{lemma}

\begin{proof}
\leanok
\uses{lem:one_le_C_MoserAnchor}

Both branches of the definition return a value $\geq C_{\mathrm{MA}}(d) \geq 1$.
The third statement follows by transitivity from the corresponding bound for
$C_{\mathrm{MA}}$.
\end{proof}

\subsection{Geometric sequences}

\begin{definition}[Nested Moser radii]
\label{def:moserRadius}
\lean{DeGiorgi.moserRadius}
\leanok
\begin{equation*}
  r_n := \frac{1 + 2^{-n}}{2}, \qquad n \geq 0.
\end{equation*}
Thus $r_0 = 1$ and $r_n \searrow 1/2$.
\end{definition}

\begin{lemma}[Properties of the radii]
\label{lem:moserRadius_zero}
\lean{DeGiorgi.moserRadius_zero, DeGiorgi.moserRadius_le_one, DeGiorgi.moserRadius_pos, DeGiorgi.moserRadius_gt_half, DeGiorgi.moserRadius_gap, DeGiorgi.moserRadius_succ_lt}
\leanok
\uses{def:moserRadius}
For all $n \in \N$:
\begin{enumerate}
  \item $r_0 = 1$ \leanname{moserRadius\_zero};
  \item $r_n \leq 1$ \leanname{moserRadius\_le\_one};
  \item $0 < r_n$ \leanname{moserRadius\_pos};
  \item $1/2 < r_n$ \leanname{moserRadius\_gt\_half};
  \item $r_n - r_{n+1} = 2^{-(n+2)}$ \leanname{moserRadius\_gap};
  \item $r_{n+1} < r_n$ \leanname{moserRadius\_succ\_lt}.
\end{enumerate}
\end{lemma}

\begin{proof}
\leanok

Each follows by direct algebra: e.g.\ for (5),
$r_n - r_{n+1} = (2^{-n} - 2^{-(n+1)})/2 = 2^{-(n+2)}$.
\end{proof}

\begin{definition}[Exponent sequence]
\label{def:moserExponentSeq}
\lean{DeGiorgi.moserExponentSeq}
\leanok
\uses{def:moserChi}
Given $p_0 \in \R$ and $n \in \N$, set
\begin{equation*}
  p_n := p_0 \chi^{n}.
\end{equation*}
\end{definition}

\begin{lemma}
\label{lem:moserExponentSeq_zero}
\lean{DeGiorgi.moserExponentSeq_zero, DeGiorgi.moserExponentSeq_succ, DeGiorgi.moserExponentSeq_pos, DeGiorgi.moserExponentSeq_ge_initial, DeGiorgi.one_lt_moserExponentSeq, DeGiorgi.inv_moserExponentSeq}
\leanok
\uses{def:moserChi, def:moserExponentSeq}
With $\chi = d/(d-2)$ and $q = \chi^{-1}$:
\begin{enumerate}
  \item $p_0 = p_0$ and $p_{n+1} = \chi p_n$ \leanname{moserExponentSeq\_zero}, \leanname{moserExponentSeq\_succ};
  \item if $d>2$ and $p_0 > 0$, then $p_n > 0$ \leanname{moserExponentSeq\_pos};
  \item if $d>2$ and $p_0 \geq 0$, then $p_0 \leq p_n$ \leanname{moserExponentSeq\_ge\_initial};
  \item if $d>2$ and $p_0 > 1$, then $p_n > 1$ \leanname{one\_lt\_moserExponentSeq};
  \item if $d>2$ and $p_0 > 0$, then $1/p_n = q^n / p_0$ \leanname{inv\_moserExponentSeq}.
\end{enumerate}
\end{lemma}

\begin{proof}
\leanok
\uses{lem:moserDecayRatio_nonneg}

All four are immediate consequences of $\chi > 1$ and $p_n = p_0 \chi^n$;
identity (5) uses $\chi^{-n} = q^n$.
\end{proof}

\subsection{Regularized power cutoffs}

For the cutoff--power test functions of the iteration we set
\begin{equation*}
  \mathrm{PC}_d(\eta,u,p)(x) := \eta(x)\,|\!\max(u(x),0)|^{p/2}
  \quad (\text{\leanname{moserPowerCutoff}}),
\end{equation*}
the smooth cutoff $\eta$ being supported in a slightly larger ball than the
inner ball where the bound is sought.

\begin{lemma}[Support property]
\label{lem:moserPowerCutoff_tsupport_subset}
\lean{DeGiorgi.moserPowerCutoff_tsupport_subset}
\leanok
If $\supp \eta \subseteq \Omega$, then
$\supp \mathrm{PC}_d(\eta,u,p) \subseteq \Omega$.
\end{lemma}

\begin{proof}
\leanok

Pointwise $\mathrm{PC}_d(\eta,u,p)(x) = 0$ whenever $\eta(x)=0$.
\end{proof}

\subsection{Sobolev preparation and the energy bound}

The hard nonlinear core of the Moser iteration is the construction of a
$W^{1,2}$ witness for the cutoff power $\mathrm{PC}_d(\eta,u,p)$ together with
the matching energy estimate.

\begin{theorem}[Cutoff--power energy bound]
\label{thm:moser-cutoff-energy}
\lean{DeGiorgi.moser_powerCutoff_memW1p_energy_of_subsolution_core}
\leanok
\uses{def:MemW1pWitness, def:ellipticityRatio, def:IsSubsolution}
Let $d > 2$, $A$ a normalized elliptic coefficient on $\Bz{1}$, $u$ a
subsolution of $A$, and $\eta \in C^{\infty}_c(\Bz{s})$ a smooth cutoff with
$0 \leq \eta \leq 1$ and $\|\nabla \eta\|_{L^\infty} \leq C_\eta$. Suppose
$0 < s \leq 1$, $1 < p$, and $\int_{\Bz{s}} |u_+|^p < \infty$. Then there
exists a $W^{1,2}$ witness for $v := \mathrm{PC}_d(\eta,u,p)$ on $\Bz{s}$
satisfying
\begin{equation*}
  \int_{\Bz{s}} \|\nabla v\|^{2}\,dx
  \;\leq\;
  2\,C_\eta^{2}\,\Big(\Lambda\,\Big(\frac{p}{p-1}\Big)^{2} + 1\Big)
  \int_{\Bz{s}} |u_+|^{p}\,dx.
\end{equation*}
\end{theorem}

\begin{proof}
\leanok
\uses{def:HasWeakPartialDeriv, lem:weak_lp_closure, def:MemW1p, def:MemW01p, lem:witness_add, lem:MemW1pWitness_restrict, lem:MemW1pWitness_smul, lem:weak_grad_norm_memLp, lem:zero_outside_tsupport, thm:memW01p_of_compactly_supported, def:CampanatoBall, def:HasCampanatoBound, thm:chain-rule-open, lem:rescale-witness, thm:w12-approx, lem:MemW1pWitness_sub_const, lem:MemW1pWitness_mul_smooth_bounded, def:EllipticCoeff, lem:EllipticCoeff_ae_coercive_nonneg, lem:EllipticCoeff_mulVec_sq_le, lem:EllipticCoeff_quadratic_upper, def:bilinFormIntegrandOfCoeff, lem:aestronglyMeasurable_bilinFormIntegrandOfCoeff, lem:bilinFormOfCoeff_eq_left, def:rescaleToUnitBall, def:rescaleCoeffToUnitBall, thm:rescaleToUnitBall_isSubsolution, def:EllipticCoeff_restrict, lem:memH01_indicator_ball_of_closedBall_subset, lem:bilinForm_ball_restrict_eq_zeroExtend, def:positivePartSub, def:positivePartSub_memW1pWitness_on_ball, cor:linfty_subsolution_DeGiorgi_exists_threshold}

The witness for $v := \mathrm{PC}_d(\eta, u, p) = \eta\,(u_+)^{p/2}$ is
constructed in \leanname{moserPowerCutoff\_memW1pWitness}, which we now
unpack. First, a $W^{1,2}$ witness for $(u_+)^{p/2}$ on $\Bz{s}$ is obtained
from the truncation $u_+$ (whose witness comes from
\leanname{positivePartSub\_memW1pWitness\_on\_ball}) composed with the
power $t \mapsto t^{p/2}$ via the $W^{1,2}$ chain rule
\leanname{sobolev\_chain\_rule\_unitBall}; the local integrability
$\int (u_+)^{p} < \infty$ together with $p > 1$ ensures the chain-rule
hypotheses (locally Lipschitz factor on $[0,\infty)$). The product
$\eta\,(u_+)^{p/2}$ is then an honest $W^{1,2}$ function via
\leanname{MemW1pWitness.mul\_smooth\_bounded\_p}, with weak gradient
$\nabla v = \eta\,(p/2)(u_+)^{p/2-1}\nabla u_+ + (u_+)^{p/2}\nabla \eta$.

For the energy bound, test the subsolution inequality against
$\varphi := \eta^{2}(u_+)^{p-1}$ (this is in $H^{1}_{0}(\Bz{s})$ thanks to
the support inclusion of $\eta$ and the integrability hypothesis on $u_+$).
The pointwise expansion of $\langle A\nabla u, \nabla \varphi\rangle$ via
\leanname{bilinFormIntegrand\_mul\_smooth\_eq} yields
\begin{multline*}
  \int_{\Bz{s}} (p-1)\,\eta^{2}\,(u_+)^{p-2}\,\langle A\nabla u_+, \nabla u_+\rangle\,dx
  \\
  \le -2\int_{\Bz{s}} \eta\,(u_+)^{p-1}\,\langle A\nabla u_+, \nabla \eta\rangle\,dx,
\end{multline*}
where on the left we discarded the supersolution-favourable term and on the
right we used the subsolution inequality. The right-hand side is bounded by
Cauchy--Schwarz against $\langle A\cdot, \cdot\rangle$ followed by AM--GM
with weight $\varepsilon = (p-1)/2$:
\begin{equation*}
  2|\eta\,(u_+)^{p-1}\langle A\nabla u_+,\nabla \eta\rangle|
  \le \frac{p-1}{2}\eta^{2}(u_+)^{p-2}\langle A\nabla u_+,\nabla u_+\rangle
  + \frac{2}{p-1}(u_+)^{p}\langle A\nabla \eta,\nabla \eta\rangle.
\end{equation*}
Absorbing the first term into the left-hand side leaves
$(p-1)/2 \cdot \int \eta^{2}(u_+)^{p-2}\langle A\nabla u_+,\nabla u_+\rangle$,
and the residual term is bounded by $\Lambda C_\eta^{2}\int (u_+)^{p}$ via
the upper ellipticity bound for $A$ and $\|\nabla \eta\| \le C_\eta$.

Finally, expanding $\nabla v = \eta\,(p/2)(u_+)^{p/2-1}\nabla u_+ +
(u_+)^{p/2}\nabla \eta$, the pointwise inequality
$\|\nabla v\|^{2} \le 2\,\eta^{2}(p/2)^{2}(u_+)^{p-2}\|\nabla u_+\|^{2}
+ 2(u_+)^{p}\|\nabla \eta\|^{2}$ together with the absorbed Caccioppoli
estimate gives
\begin{equation*}
  \int_{\Bz{s}} \|\nabla v\|^{2}\,dx
   \le 2\,C_\eta^{2}\Bigl(\Lambda\bigl(\tfrac{p}{p-1}\bigr)^{2} + 1\Bigr)
   \int_{\Bz{s}} (u_+)^{p}\,dx,
\end{equation*}
the displayed energy inequality.
\end{proof}

\begin{theorem}[Zero trace upgrade]
\label{thm:moser_powerCutoff_memW01p_energy_of_subsolution}
\lean{DeGiorgi.moser_powerCutoff_memW01p_energy_of_subsolution}
\leanok
\uses{def:MemW1pWitness, def:ellipticityRatio, def:IsSubsolution}
Under the hypotheses of Theorem~\ref{thm:moser-cutoff-energy}, the witness
moreover satisfies $v \in W^{1,2}_{0}(\Bz{s})$.
\end{theorem}

\begin{proof}
\leanok
\uses{def:MemW1p, def:MemW01p, thm:memW01p_of_compactly_supported, lem:moserPowerCutoff_tsupport_subset, thm:moser-cutoff-energy}

Since $\supp \eta \subseteq \Bz{s}$ is compact, $\supp v \subseteq \Bz{s}$ is
likewise compactly contained. The criterion
\leanname{memW01p\_of\_memW1p\_of\_tsupport\_subset} then upgrades $v$ from
$W^{1,p}$ to $W^{1,p}_0$.
\end{proof}

\subsection{One-step Moser pre-estimate}

\begin{theorem}[Concentric ball pre-estimate]
\label{thm:moser-preMoser}
\lean{DeGiorgi.moser_preMoser}
\leanok
\uses{def:ellipticityRatio, def:IsSubsolution, def:C_MoserAnchor, def:moserChi}
Let $d > 2$, $A$ as above, $u$ a subsolution, and $0 < r < s \leq 1$.
Suppose $1 < p$ and $\int_{\Bz{s}} |u_+|^p < \infty$. Then $|u_+|^{\chi p}$
is integrable on $\Bz{r}$ and
\begin{multline*}
  \Big(\int_{\Bz{r}} |u_+|^{\chi p}\,dx\Big)^{1/(\chi p)}
  \\ \leq\;
  \Big(\frac{C_{\mathrm{MA}}(d)}{(s-r)^{2}}\,
    \Big(\Lambda\Big(\frac{p}{p-1}\Big)^{2} + 1\Big)\Big)^{1/p}
  \,\Big(\int_{\Bz{s}} |u_+|^{p}\,dx\Big)^{1/p}.
\end{multline*}
\end{theorem}

\begin{proof}
\leanok
\uses{def:MemW1pWitness, def:MemW01p, lem:zero_outside_tsupport, thm:zero_extend, thm:weak_grad_unique, def:C_gns, cor:sobolev_W01p, def:CampanatoBall, def:HasCampanatoBound, lem:smoothTransition_nnnorm_deriv_bounded, def:Cutoff, thm:Cutoff_exists, lem:one_le_C_MoserAnchor, lem:moserDecayRatio_nonneg, lem:moserPowerCutoff_tsupport_subset, thm:moser-cutoff-energy, thm:moser_powerCutoff_memW01p_energy_of_subsolution}

Choose the midpoint radius $R := (r+s)/2 \in (r,s)$ and a smooth radial cutoff
$\eta$ with $\eta = 1$ on $\Bz{r}$, $\supp \eta \subseteq \overline{\Bz{R}}
\subseteq \Bz{s}$, $0 \leq \eta \leq 1$, and (by
\leanname{Cutoff.exists}) $\|\nabla \eta\|_{L^\infty} \leq C_\eta := 2\Mst/(s-r)$.

Set $v := \mathrm{PC}_d(\eta,u,p) = \eta\,(u_+)^{p/2}$. By
Theorem~\ref{thm:moser-cutoff-energy} we get $v \in W^{1,2}_0(\Bz{s})$ with
\begin{equation*}
  \int_{\Bz{s}} \|\nabla v\|^{2}\,dx
  \;\leq\;
  2\,C_\eta^{2}\,\Big(\Lambda\Big(\frac{p}{p-1}\Big)^{2}+1\Big)
  \int_{\Bz{s}} (u_+)^{p}\,dx.
\end{equation*}
The Sobolev inequality \leanname{sobolev\_prepare\_on\_ball} on $\Bz{s}$ in
exponent $2 \to 2^{*} = 2d/(d-2)$ gives, with $C_{\mathrm{gns}} = C_{\mathrm{gns}}(d,2)$,
\begin{equation*}
  \Big(\int_{\Bz{s}} v^{2^{*}}\,dx\Big)^{1/2^{*}}
  \;\leq\;
  C_{\mathrm{gns}}\,
  \Big(\int_{\Bz{s}} \|\nabla v\|^{2}\,dx\Big)^{1/2}.
\end{equation*}
On $\Bz{r}$ we have $\eta = 1$, so $v^{2^{*}} = (u_+)^{p\,d/(d-2)} = (u_+)^{\chi p}$.
Combining with the energy bound and using $C_\eta^2 = 4\Mst^2/(s-r)^2$ together
with $8\Mst^{2}\,C_{\mathrm{gns}}^{2} \leq C_{\mathrm{MA}}(d)$
(\leanname{cutoff\_sobolev\_anchor\_le\_C\_MoserAnchor}) yields
\begin{equation*}
  \Big(\int_{\Bz{r}} (u_+)^{\chi p}\,dx\Big)^{2/2^{*}}
  \;\leq\;
  \frac{C_{\mathrm{MA}}(d)}{(s-r)^{2}}
  \Big(\Lambda\Big(\frac{p}{p-1}\Big)^{2}+1\Big)
  \int_{\Bz{s}} (u_+)^{p}\,dx.
\end{equation*}
Raising to the power $1/p$ and using $2/(2^{*}\,p)= 1/(\chi p)$ gives the
claim.
\end{proof}

\subsection{Iteration}

We now iterate the pre-estimate along $(r_n,p_n)$.

\begin{definition}[Iteration quantities]
\label{def:moserIterIntegral}
\lean{DeGiorgi.moserIterIntegral, DeGiorgi.moserIterNorm, DeGiorgi.moserStepConst, DeGiorgi.moserLinftyBoundPow, DeGiorgi.moserLinftyMajorant}
\leanok
\uses{def:ellipticityRatio, def:C_MoserAnchor, def:C_Moser, def:moserRadius, def:moserExponentSeq}
Set
\begin{align*}
  I_n(u) &:= \int_{\Bz{r_n}} (u_+)^{p_n}\,dx
    && (\text{\leanname{moserIterIntegral}}), \\
  N_n(u) &:= I_n(u)^{1/p_n}
    && (\text{\leanname{moserIterNorm}}), \\
  \mathrm{S}(A,p_0) &:= 32\,C_{\mathrm{MA}}(d)\,\Lambda\,\Big(\frac{p_0}{p_0-1}\Big)^{2}
    && (\text{\leanname{moserStepConst}}).
\end{align*}
The product majorant produced by iteration is, for $n \geq 0$,
\begin{equation*}
  M_n(A,u,p_0) :=
  \Big(
    \mathrm{S}(A,p_0)^{\sum_{i=0}^{n-1} q^{i}}\,
    4^{\sum_{i=0}^{n-1} i\,q^{i}}\,
    \int_{\Bz{1}} (u_+)^{p_0}\,dx
  \Big)^{1/p_0}.
\end{equation*}
We further define the target majorant
\begin{align*}
  B_\infty(A,u,p_0) &:= C_{\mathrm{Moser}}(d)\,\Lambda^{d/2}\,
    \Big(\frac{p_0}{p_0-1}\Big)^{d}\,\int_{\Bz{1}} (u_+)^{p_0}\,dx
    && (\text{\leanname{moserLinftyBoundPow}}), \\
  M_\infty(A,u,p_0) &:= B_\infty(A,u,p_0)^{1/p_0}
    && (\text{\leanname{moserLinftyMajorant}}).
\end{align*}
\end{definition}

\begin{theorem}[Step--majorant inequality]
\label{thm:moser-step-maj}
\lean{DeGiorgi.moser_step_majorant_le}
\leanok
\uses{def:ellipticityRatio, def:C_MoserAnchor, def:moserChi, def:moserRadius, def:moserExponentSeq, def:moserIterIntegral}
For $d > 2$, $1 < p_0$ and any $n \geq 0$,
\begin{multline*}
  \Big(\frac{C_{\mathrm{MA}}(d)}{(r_n-r_{n+1})^{2}}\,
    \Big(\Lambda\Big(\frac{p_n}{p_n-1}\Big)^{2}+1\Big)\Big)^{1/p_n}
  M_n(A,u,p_0) \\
  \leq\;
  M_{n+1}(A,u,p_0).
\end{multline*}
\end{theorem}

\begin{proof}
\leanok
\uses{lem:one_le_C_MoserAnchor, lem:moserDecayRatio_nonneg, lem:moserRadius_zero, lem:moserExponentSeq_zero}

First, $r_n - r_{n+1} = 2^{-(n+2)}$ so
$C_{\mathrm{MA}}/(r_n-r_{n+1})^2 = C_{\mathrm{MA}}\,4^{n+2} = 16\,C_{\mathrm{MA}}\,4^{n}$.
Second, $p_n/(p_n-1) \leq p_0/(p_0-1)$ since $t \mapsto t/(t-1)$ is decreasing
on $(1,\infty)$, hence
\begin{equation*}
  \Lambda\Big(\frac{p_n}{p_n-1}\Big)^{2}+1
  \;\leq\; 2\,\Lambda\Big(\frac{p_0}{p_0-1}\Big)^{2}.
\end{equation*}
Combining gives
\begin{equation*}
  \frac{C_{\mathrm{MA}}}{(r_n-r_{n+1})^{2}}\,
    \Big(\Lambda\Big(\frac{p_n}{p_n-1}\Big)^{2}+1\Big)
  \;\leq\;
  \mathrm{S}(A,p_0)\,4^{n}.
\end{equation*}
Raising to $1/p_n = q^n/p_0$ and multiplying by $M_n$ produces the inner factor
$\mathrm{S}^{q^n/p_0} 4^{n q^n/p_0}$, which precisely promotes the partial sums
in the exponents from $\sum_{i<n}$ to $\sum_{i\leq n}$. The resulting expression
equals $M_{n+1}(A,u,p_0)$ by the definition of $M_{n+1}$.
\end{proof}

\begin{theorem}[Iteration bound]
\label{thm:moser-iter-bound}
\lean{DeGiorgi.moser_iteration_bound}
\leanok
\uses{def:ellipticityRatio, def:IsSubsolution, def:moserRadius, def:moserExponentSeq, def:moserIterIntegral}
For $d > 2$, $1 < p_0$, $u$ a subsolution of $A$, and
$\int_{\Bz{1}} |u_+|^{p_0} < \infty$, for every $n \in \N$ we have
$|u_+|^{p_n}$ integrable on $\Bz{r_n}$ and
\begin{equation*}
  N_n(u) \;\leq\; M_n(A,u,p_0).
\end{equation*}
\end{theorem}

\begin{proof}
\leanok
\uses{def:C_MoserAnchor, lem:one_le_C_MoserAnchor, def:moserChi, lem:moserRadius_zero, lem:moserExponentSeq_zero, thm:moser-preMoser, thm:moser-step-maj}

By induction on $n$. The case $n=0$ uses $r_0=1$, $p_0 = p_0$, and
$M_0(A,u,p_0)= (\int_{\Bz{1}} (u_+)^{p_0})^{1/p_0}$ by definition. For the
inductive step, the pre-estimate Theorem~\ref{thm:moser-preMoser} applied with
$(p,r,s) = (p_n,r_{n+1},r_n)$ gives integrability of $|u_+|^{\chi p_n}$ on
$\Bz{r_{n+1}}$ and
\begin{multline*}
  N_{n+1}(u)
  \;\leq\;
  \Big(\frac{C_{\mathrm{MA}}}{(r_n-r_{n+1})^{2}}
    \Big(\Lambda\Big(\frac{p_n}{p_n-1}\Big)^{2}+1\Big)\Big)^{1/p_n}
  N_n(u).
\end{multline*}
The inductive hypothesis $N_n(u) \leq M_n(A,u,p_0)$ and
Theorem~\ref{thm:moser-step-maj} then yield $N_{n+1}(u) \leq M_{n+1}(A,u,p_0)$.
\end{proof}

\begin{lemma}[Geometric majorant absorption]
\label{lem:moser-geom-maj}
\lean{DeGiorgi.moser_geometric_majorant}
\leanok
\uses{def:ellipticityRatio, def:moserChi, def:moserIterIntegral}
For $d > 2$, $1 < p_0$, and all $n$,
\begin{equation*}
  M_n(A,u,p_0) \;\leq\; M_\infty(A,u,p_0).
\end{equation*}
\end{lemma}

\begin{proof}
\leanok
\uses{def:C_MoserAnchor, lem:one_le_C_MoserAnchor, lem:moserDecayRatio_nonneg, def:C_Moser}

Set $S := \sum_{i=0}^{n-1} q^{i}$, $T := \sum_{i=0}^{n-1} i q^{i}$,
$B := 32\,C_{\mathrm{MA}}(d)$, $r := p_0/(p_0-1)$. Then:
$S \leq d/2$ \leanname{moserDecayRatio\_sum\_le\_half\_dim}, since
$\sum_{i\geq 0} q^{i} = 1/(1-q) = d/2$; also $T \leq \sum_{i\geq 0} i\,q^{i}$
\leanname{moserDecayRatio\_weighted\_sum\_le\_tsum} by summability.
Now $\mathrm{S}(A,p_0) = B \Lambda r^{2}$, so
\begin{equation*}
  \mathrm{S}(A,p_0)^{S}
  = B^{S} \Lambda^{S} (r^{2})^{S}.
\end{equation*}
Using $1 \leq B$, $1 \leq \Lambda$, and $1 \leq r$, monotonicity of $t \mapsto
t^{e}$ in the exponent gives $B^{S} \leq B^{d/2}$, $\Lambda^{S} \leq
\Lambda^{d/2}$, and $(r^{2})^{S} \leq r^{d}$ (because $2S \leq d$). Also $4^{T}
\leq 4^{\sum_{i\geq 0} i q^{i}}$. Multiplying these inequalities and using the
explicit definition of $C_{\mathrm{Moser}}(d)$ in the case $d>2$:
\begin{equation*}
  B^{d/2}\,4^{\sum_{i\geq 0} i q^{i}} \;\leq\; C_{\mathrm{Moser}}(d),
\end{equation*}
we obtain
\begin{equation*}
  \mathrm{S}(A,p_0)^{S}\,4^{T}\,I_0(u)
  \;\leq\;
  C_{\mathrm{Moser}}(d)\,\Lambda^{d/2}\,r^{d}\,I_0(u)
  = B_\infty(A,u,p_0).
\end{equation*}
Raising to the power $1/p_0$ yields $M_n \leq M_\infty$.
\end{proof}

\subsection{Closeout to $L^\infty$}

\begin{lemma}[Closeout from iterated bounds]
\label{lem:moser-closeout}
\leanok
\uses{def:ellipticityRatio, def:moserRadius, def:moserExponentSeq, def:moserIterIntegral}
\noindent\textit{(Formalized as private declaration \leanname{DeGiorgi.moser\_ae\_closeout}.)}\par
Fix $d > 2$, $1 < p_0$, and assume that for every $n \in \N$,
$|u_+|^{p_n}$ is integrable on $\Bz{r_n}$ and $N_n(u) \leq M_\infty(A,u,p_0)$.
Then for almost every $x \in \Bz{1/2}$,
\begin{equation*}
  |u_+(x)|^{p_0} \;\leq\; B_\infty(A,u,p_0).
\end{equation*}
\end{lemma}

\begin{proof}
\leanok
\uses{def:moserChi, lem:moserDecayRatio_nonneg, lem:C_MoserAnchor_le_C_Moser, lem:moserRadius_zero, lem:moserExponentSeq_zero}

Set $\mu := \mathrm{vol}\,|\,\Bz{1/2}$, $f(x) := |u_+(x)|$, $g(x) := f(x)^{p_0/2}$,
and $K := B_\infty(A,u,p_0)^{1/2}$. Note $K \geq 0$.
For each $n$, write $q_n := 2 \chi^n$. We claim that for every $n$
and every $c > K$,
\begin{equation*}
  \mu_{\mathrm{r}}\{x : c \leq g(x)\} \;\leq\; (K/c)^{q_n}.
\end{equation*}
Indeed, on $\{c \leq g\}$ we have $c^{q_n} \leq g^{q_n}$, and Markov's
inequality gives
\begin{equation*}
  c^{q_n}\,\mu_{\mathrm{r}}\{c \leq g\}
  \;\leq\;
  \int g^{q_n}\,d\mu
  = \int_{\Bz{1/2}} f^{p_n}\,dx
  \;\leq\;
  \int_{\Bz{r_n}} f^{p_n}\,dx = I_n(u),
\end{equation*}
where we used $g^{q_n} = f^{p_0 q_n / 2} = f^{p_n}$ and
$\Bz{1/2} \subseteq \Bz{r_n}$. By hypothesis $I_n(u)^{1/p_n} = N_n(u) \leq
M_\infty$, so raising to $p_n$ and using $1/q_n = (1/p_n)(p_0/2)$ we get
$I_n(u) \leq M_\infty^{q_n p_0/2 \cdot 2/p_0} = M_\infty^{q_n p_0/p_0}$\,;
more cleanly, $I_n^{1/q_n} \leq M_\infty^{p_0/2} = K$. Hence
$\mu_{\mathrm{r}}\{c\leq g\} \leq (K/c)^{q_n}$.

As $\chi > 1$, $q_n \to \infty$, so $(K/c)^{q_n} \to 0$ for every $c > K$.
Therefore $\mu_{\mathrm{r}}\{c \leq g\} = 0$ for all $c > K$, and a countable
union over $c = K + 1/(m+1)$ shows $g \leq K$ a.e.\ on $\Bz{1/2}$. Squaring,
$f^{p_0} = g^{2} \leq K^{2} = B_\infty$.
\end{proof}

\begin{theorem}[Moser $L^p \to L^\infty$ for subsolutions]
\label{thm:linfty-Moser}
\lean{DeGiorgi.linfty_subsolution_Moser}
\leanok
\uses{def:ellipticityRatio, def:IsSubsolution, def:C_Moser}
Let $d > 2$, $A$ a normalized elliptic coefficient on $\Bz{1}$, $u$ a
subsolution, and $1 < p_0$ with $\int_{\Bz{1}} |u_+|^{p_0} < \infty$. Then for
almost every $x \in \Bz{1/2}$,
\begin{equation*}
  |u_+(x)|^{p_0}
  \;\leq\;
  C_{\mathrm{Moser}}(d)\,\Lambda^{d/2}\,
  \Big(\frac{p_0}{p_0-1}\Big)^{d}
  \int_{\Bz{1}} |u_+|^{p_0}\,dx.
\end{equation*}
\end{theorem}

\begin{proof}
\leanok
\uses{def:moserRadius, def:moserExponentSeq, thm:moser-preMoser, def:moserIterIntegral, thm:moser-iter-bound, lem:moser-geom-maj, lem:moser-closeout}

The headline Moser bound is the composition of three pieces, all assembled
along the discrete sequence of radii $r_n$ and exponents $p_n$.

\emph{Step 1: iterated step bound.} Theorem~\ref{thm:moser-iter-bound}
(\leanname{moser\_iteration\_bound}) is applied as a single induction on
$n$: at each step, the concentric-ball pre-estimate
Theorem~\ref{thm:moser-preMoser} on $(r_{n+1}, r_n)$ with exponent $p_n$
upgrades $L^{p_n}$ control of $u_+$ on $\Bz{r_n}$ to $L^{p_{n+1}} = L^{\chi p_n}$
control on $\Bz{r_{n+1}}$, with multiplicative constant the step constant
$\mathrm{S}_n$ depending on $\Lambda$, $C_{\mathrm{MA}}(d)$, the dimension,
and the gap $r_n - r_{n+1} = 2^{-(n+2)}$. Telescoping over $i < n$ produces
the bound
\begin{equation*}
  N_n(u)
   \le \Bigl(\prod_{i=0}^{n-1} \mathrm{S}_i\Bigr) \cdot N_0(u)
   = M_n(A, u, p_0).
\end{equation*}
At every stage the integrability of $|u_+|^{p_n}$ on $\Bz{r_n}$ is propagated
forward by the same induction, providing the second half of the conclusion.

\emph{Step 2: geometric absorption of the product.}
Lemma~\ref{lem:moser-geom-maj} (\leanname{moser\_geometric\_majorant})
bounds the infinite product $\prod_{i \ge 0} \mathrm{S}_i$ by the closed-form
constant $M_\infty(A,u,p_0)$. The key two summation identities are
$\sum_{i \ge 0} q^{i} \le d/2$
(\leanname{moserDecayRatio\_sum\_le\_half\_dim}) and the weighted variant
$\sum_{i \ge 0} i\,q^{i}$
(\leanname{moserDecayRatio\_weighted\_sum\_le\_tsum}); these are precisely
what is needed to convert the product of $\mathrm{S}_i$, after taking
logarithms and using $1/p_i = q^{i}/p_0$, into a finite expression involving
$C_{\mathrm{Moser}}(d)$, $\Lambda^{d/2}$ and the explicit
$\bigl(p_0/(p_0-1)\bigr)^{d}$ factor visible in the conclusion. Composing with
Step 1 gives $N_n(u) \le M_\infty(A, u, p_0)$ for every $n$.

\emph{Step 3: a.e.\ closeout via Markov.}
Lemma~\ref{lem:moser-closeout} (\leanname{moser\_ae\_closeout}) converts the
uniform $L^{p_n}$ bound at every stage into an a.e.\ pointwise bound on the
intersection $\Bz{1/2} \subseteq \bigcap_n \Bz{r_n}$. The argument is a
Markov-type estimate: for every $c > K := M_\infty(A,u,p_0)^{1/2}$, the
super-level set $\{c \le f\}$ on $\Bz{1/2}$ has measure at most
$(K/c)^{q_n}$, and since $q_n \to \infty$ this measure is zero. Taking a
countable union over $c = K + 1/(m+1)$ shows $f \le K$ a.e.\ on $\Bz{1/2}$,
and squaring gives $|u_+|^{p_0} \le K^{2}$ a.e.\ on $\Bz{1/2}$, which is the
displayed Moser bound.
\end{proof}

\begin{remark}[$p_0 = 2$ anchor, \leanname{linfty\_subsolution\_Moser\_two}]
The case $p_0 = 2$ already follows from the normalized De Giorgi $L^2 \to
L^\infty$ estimate of Section 5: with $c := C_{\mathrm{DG}}^{\mathrm{sub}}(d)
\Lambda^{d/4}$, the De Giorgi bound gives
$\max(u(x),0) \leq c\,(\int_{\Bz{1}} |u_+|^{2}\,dx)^{1/2}$ a.e.\ on
$\Bz{1/2}$. Squaring and using
$C_{\mathrm{DG}}^{\mathrm{sub}}(d)^{2} \leq C_{\mathrm{Moser}}(d)$ together with
$2/(2-1) = 2$ recovers the $p_0=2$ case of Theorem~\ref{thm:linfty-Moser}.
\end{remark}

%% --- Section 8 ---

\section{Supersolution Estimates}

We now turn to positive supersolutions of the same operator. The main outputs
of this section are the two stages of the weak Harnack inequality:
\emph{stage one (inverse)} produces an a.e.\ lower bound on $u$ from inverse
$L^{p}$ control of $u$, and \emph{stage two (forward)} produces a low-power
$L^{p}$-on-$L^{p}$ Caccioppoli-type bound. The architecture mirrors Section 7,
with the role of $u_+$ replaced by $u^{-1}$ (resp.\ $u$ at low powers).

Throughout this section, $u : E \to \R$ is assumed to be strictly positive on
$\Bz{1}$ and a supersolution of $A$. We write $\mu_{1/2}$ for
$\mathrm{vol}\,|\,\Bz{1/2}$.

\subsection{The weighted Caccioppoli inequality}

The abstract energy tool driving every iteration in this section is the
following weighted Caccioppoli inequality.

\begin{theorem}[Weighted Caccioppoli]
\label{thm:weighted-cacc}
\lean{DeGiorgi.weighted_caccioppoli_of_supersolution}
\leanok
\uses{def:MemW1pWitness, def:MemW01p, def:ellipticityRatio, def:bilinFormIntegrandOfCoeff, def:IsSubsolution}
Let $\Omega \subseteq E$ be open, $A$ a normalized elliptic coefficient on
$\Omega$, $u$ a supersolution with $u \in W^{1,2}(\Omega)$, and let
$\Psi : \R \to \R$ be smooth with $\Psi(0)=0$, $|\Psi'| \leq M$, $\Psi(u) \geq 0$
on $\Omega$, and $\Psi'(u) \leq 0$ on $\Omega$ (with $\Psi(u)=0$ wherever
$\Psi'(u)=0$). Let $\varphi \in C^{\infty}_c(\Omega)$ satisfy $|\varphi| \leq 1$
and $\|\nabla \varphi\|_{L^\infty} \leq C_\varphi$, with the integrability
\begin{equation*}
  \int_{\Omega} \|\nabla \varphi\|^{2}\,
  \frac{\Psi(u)^{2}}{-\Psi'(u)}\,dx < \infty
\end{equation*}
(interpreted as $0$ where $\Psi'(u)=0$). Then
\begin{multline*}
  \int_{\Omega} \varphi^{2}\,(-\Psi'(u))\,\langle A\nabla u,\nabla u\rangle\,dx
  \\ \leq\;
  4\,\Lambda \int_{\Omega} \|\nabla \varphi\|^{2}\,
  \frac{\Psi(u)^{2}}{-\Psi'(u)}\,dx.
\end{multline*}
\end{theorem}

\begin{proof}
\leanok
\uses{def:HasWeakPartialDeriv, prop:witness_mul_smooth_bounded, thm:chain-rule-open, lem:EllipticCoeff_ae_coercive_nonneg, lem:EllipticCoeff_mulVec_sq_le, lem:aestronglyMeasurable_bilinFormIntegrandOfCoeff, lem:bilinFormIntegrand_mul_smooth_eq, lem:integrable_bounded_mul_bilinFormIntegrand}

By the supersolution property tested against $\varphi^{2}\Psi(u) \geq 0$,
\begin{equation*}
  0 \leq \int_{\Omega} \langle A\nabla u, \nabla(\varphi^{2}\Psi(u))\rangle\,dx.
\end{equation*}
Using the chain rule for $W^{1,2}$ functions, $\nabla(\Psi(u)) = \Psi'(u)\nabla u$,
and the product rule $\nabla(\varphi^{2}\Psi(u)) = \varphi^{2}\Psi'(u)\nabla u
+ 2\varphi\Psi(u)\nabla \varphi$. The bilinear form expansion
(\leanname{bilinFormIntegrand\_mul\_smooth\_eq}) yields the pointwise identity
\begin{equation*}
  \langle A\nabla u, \nabla(\varphi^{2}\Psi(u))\rangle
  = \varphi^{2}\Psi'(u)\,\langle A\nabla u,\nabla u\rangle
  + 2\varphi\Psi(u)\,\langle A\nabla u,\nabla \varphi\rangle.
\end{equation*}
Since $\Psi'(u) \leq 0$, the principal term is non-positive; rearranging
\begin{equation*}
  \int \varphi^{2}\,(-\Psi'(u))\,\langle A\nabla u,\nabla u\rangle\,dx
  \;\leq\;
  2 \int \varphi\,\Psi(u)\,\langle A\nabla u,\nabla \varphi\rangle\,dx.
\end{equation*}
By Cauchy--Schwarz with respect to the symmetric form $\langle A\cdot,\cdot\rangle$
and AM--GM with weight $\varepsilon = 1/2$,
\begin{multline*}
  2|\varphi\Psi(u)\,\langle A\nabla u,\nabla \varphi\rangle|
  \\ \leq\;
  \tfrac12\,\varphi^{2}(-\Psi'(u))\langle A\nabla u,\nabla u\rangle
  + 2\,\frac{\Psi(u)^{2}}{-\Psi'(u)}\,\langle A\nabla \varphi,\nabla \varphi\rangle.
\end{multline*}
Absorbing the first term into the left-hand side and bounding
$\langle A\nabla \varphi,\nabla \varphi\rangle \leq \Lambda \|\nabla \varphi\|^{2}$
(ellipticity \leanname{NormalizedEllipticCoeff.upper}) yields the stated inequality
with constant $4\Lambda$.
\end{proof}

\subsection{Reverse Holder inequalities (one--step)}

\paragraph{Inverse case.}

\begin{definition}[Inverse-power cutoff]
\label{def:superPowerCutoff}
\lean{DeGiorgi.superPowerCutoff}
\leanok
For a smooth $\eta$, positive $u$, and $p > 0$,
\begin{equation*}
  \mathrm{PCI}_d(\eta,u,p)(x) := \eta(x)\,|u(x)^{-1}|^{p/2}.
\end{equation*}
\end{definition}

\begin{theorem}[Inverse-power one-step gain]
\label{thm:preMoser-inv}
\lean{DeGiorgi.supersolution_preMoser_inverse}
\leanok
\uses{def:ellipticityRatio, def:IsSubsolution, def:C_MoserAnchor, def:moserChi}
Let $d>2$, $A$ as above, $u > 0$ a supersolution, $p > 0$, and
$0 < r < s \leq 1$ with $\int_{\Bz{s}} |u^{-1}|^{p}\,dx < \infty$. Then
$|u^{-1}|^{\chi p}$ is integrable on $\Bz{r}$ and
\begin{multline*}
  \Big(\int_{\Bz{r}} |u^{-1}|^{\chi p}\,dx\Big)^{1/(\chi p)}
  \\ \leq\;
  \Big(\frac{C_{\mathrm{MA}}(d)}{(s-r)^{2}}\,
    \Big(\Lambda\Big(\frac{p}{1+p}\Big)^{2}+1\Big)\Big)^{1/p}
  \Big(\int_{\Bz{s}} |u^{-1}|^{p}\,dx\Big)^{1/p}.
\end{multline*}
\end{theorem}

\begin{proof}
\leanok
\uses{def:HasWeakPartialDeriv, lem:weak_lp_closure, def:MemW1p, def:MemW1pWitness, def:MemW01p, lem:witness_add, lem:MemW1pWitness_restrict, lem:MemW1pWitness_smul, lem:witness_smooth_cpt, lem:weak_grad_norm_memLp, lem:zero_outside_tsupport, thm:zero_extend, thm:weak_grad_unique, thm:memW01p_of_compactly_supported, def:C_gns, cor:sobolev_W01p, def:CampanatoBall, def:HasCampanatoBound, thm:chain-rule-open, lem:smoothTransition_nnnorm_deriv_bounded, def:Cutoff, thm:Cutoff_exists, lem:MemW1pWitness_sub_const, lem:MemW1pWitness_mul_smooth_bounded, def:EllipticCoeff, lem:EllipticCoeff_ae_coercive_nonneg, lem:EllipticCoeff_mulVec_sq_le, lem:EllipticCoeff_quadratic_upper, def:bilinFormIntegrandOfCoeff, lem:aestronglyMeasurable_bilinFormIntegrandOfCoeff, lem:one_le_C_MoserAnchor, lem:moserDecayRatio_nonneg, thm:moser-preMoser, thm:weighted-cacc, def:superPowerCutoff}

Apply Theorem~\ref{thm:weighted-cacc} with $\Psi(t) := t^{-p}$ on $(0,\infty)$
(suitably truncated), so that $\Psi'(t) = -p\,t^{-p-1} \leq 0$ and
$\Psi(t)^{2}/(-\Psi'(t)) = t^{-p-1}/p$. Choose a smooth cutoff $\eta$ with
$\eta = 1$ on $\Bz{r}$, $\supp \eta \subseteq \Bz{R}$ where $R := (r+s)/2$,
$0 \leq \eta \leq 1$, $\|\nabla \eta\|_{L^\infty} \leq C_\eta := 2\Mst/(s-r)$.
The weighted Caccioppoli inequality applied to $\varphi := \eta$ then yields
\begin{equation*}
  \int_{\Bz{s}} \eta^{2}\,p\,u^{-p-1}\,\langle A\nabla u,\nabla u\rangle\,dx
  \;\leq\;
  \frac{4\Lambda}{p}
  \int_{\Bz{s}} \|\nabla \eta\|^{2}\,u^{-p-1}\,dx.
\end{equation*}
Combined with the chain-rule identity
$\nabla(\eta\,u^{-p/2}) = (\nabla \eta) u^{-p/2} - (p/2) \eta\,u^{-p/2-1}\nabla u$
and the elementary $|a-b|^{2} \leq 2(a^{2}+b^{2})$ this gives, for the witness
$v := \mathrm{PCI}_d(\eta,u,p)$, the energy bound
\begin{equation*}
  \int_{\Bz{s}} \|\nabla v\|^{2}\,dx
  \;\leq\;
  2\,C_\eta^{2}\,\Big(\Lambda\Big(\frac{p}{1+p}\Big)^{2}+1\Big)
  \int_{\Bz{s}} |u^{-1}|^{p}\,dx;
\end{equation*}
this is the content of \leanname{superPowerCutoff\_energy\_bound}. Now apply
the Sobolev inequality \leanname{sobolev\_prepare\_on\_ball} to $v$ on $\Bz{s}$
and use $\eta = 1$ on $\Bz{r}$, exactly as in the proof of Theorem~\ref{thm:moser-preMoser}.
The constant $8\Mst^{2}\,C_{\mathrm{gns}}^{2} \leq C_{\mathrm{MA}}$ absorbs the
gradient bound, yielding the displayed reverse Holder inequality.
\end{proof}

\paragraph{Forward (low-power) case.}

\begin{definition}[Forward low-power cutoff]
\label{def:superPowerCutoffFwd}
\lean{DeGiorgi.superPowerCutoffFwd}
\leanok
For smooth $\eta$, $u > 0$, and $p \in (0,1)$,
\begin{equation*}
  \mathrm{PCF}_d(\eta,u,p)(x) := \eta(x)\,|u(x)|^{p/2}.
\end{equation*}
\end{definition}

\begin{theorem}[Forward low-power one-step gain]
\label{thm:preMoser-fwd}
\lean{DeGiorgi.supersolution_preMoser_forward}
\leanok
\uses{def:ellipticityRatio, def:IsSubsolution, def:C_MoserAnchor, def:moserChi}
Let $d>2$, $A$ as above, $u>0$ a supersolution, $0 < p < 1$, and
$0 < r < s \leq 1$ with $\int_{\Bz{s}} |u|^{p}\,dx < \infty$. Then
$|u|^{\chi p}$ is integrable on $\Bz{r}$ and
\begin{multline*}
  \Big(\int_{\Bz{r}} |u|^{\chi p}\,dx\Big)^{1/(\chi p)}
  \\ \leq\;
  \Big(\frac{C_{\mathrm{MA}}(d)}{(s-r)^{2}}\,
    \Big(\Lambda\Big(\frac{p}{1-p}\Big)^{2}+1\Big)\Big)^{1/p}
  \Big(\int_{\Bz{s}} |u|^{p}\,dx\Big)^{1/p}.
\end{multline*}
\end{theorem}

\begin{proof}
\leanok
\uses{def:HasWeakPartialDeriv, lem:weak_lp_closure, def:MemW1p, def:MemW1pWitness, def:MemW01p, lem:witness_add, lem:MemW1pWitness_restrict, lem:MemW1pWitness_smul, lem:witness_smooth_cpt, lem:weak_grad_norm_memLp, lem:zero_outside_tsupport, thm:zero_extend, thm:weak_grad_unique, thm:memW01p_of_compactly_supported, def:C_gns, cor:sobolev_W01p, def:CampanatoBall, def:HasCampanatoBound, thm:chain-rule-open, lem:smoothTransition_nnnorm_deriv_bounded, def:Cutoff, thm:Cutoff_exists, lem:MemW1pWitness_sub_const, lem:MemW1pWitness_mul_smooth_bounded, def:EllipticCoeff, lem:EllipticCoeff_ae_coercive_nonneg, lem:EllipticCoeff_mulVec_sq_le, lem:EllipticCoeff_quadratic_upper, def:bilinFormIntegrandOfCoeff, lem:aestronglyMeasurable_bilinFormIntegrandOfCoeff, lem:one_le_C_MoserAnchor, lem:moserDecayRatio_nonneg, def:superPowerCutoffFwd}

This time use $\Psi(t) := -t^{p}$ on $(0,\infty)$, so $\Psi'(t) = -p\,t^{p-1}
\leq 0$ and $\Psi(t)^{2}/(-\Psi'(t)) = t^{p+1}/p$. (Note that for $p \in (0,1)$
the function $t \mapsto t^{p}$ is concave with $\Psi'(0)=+\infty$, but
truncation issues are absorbed in \leanname{superPowerCutoffFwd\_energy\_bound}.)
The weighted Caccioppoli inequality with $\varphi := \eta$ supplies the
analogous energy estimate
\begin{equation*}
  \int_{\Bz{s}} \|\nabla v\|^{2}\,dx
  \;\leq\;
  2 C_\eta^{2}\,\Big(\Lambda\Big(\frac{p}{1-p}\Big)^{2}+1\Big)
  \int_{\Bz{s}} |u|^{p}\,dx
\end{equation*}
for $v := \mathrm{PCF}_d(\eta,u,p)$. Sobolev embedding then concludes as in the
inverse case.
\end{proof}

\subsection{Stage one (inverse iteration)}

We collect the inverse-power constants used in the iteration.

\begin{definition}[Inverse iteration quantities]
\label{def:superIterIntegralInv}
\lean{DeGiorgi.superIterIntegralInv, DeGiorgi.superIterNormInv, DeGiorgi.superStepConstInv}
\leanok
\uses{def:ellipticityRatio, def:C_MoserAnchor, def:moserRadius, def:moserExponentSeq}
For $u>0$, $p_0 > 0$, $n \in \N$, set
\begin{align*}
  I_n^{\mathrm{inv}}(u) &:= \int_{\Bz{r_n}} |u^{-1}|^{p_n}\,dx, \\
  N_n^{\mathrm{inv}}(u) &:= \big(I_n^{\mathrm{inv}}(u)\big)^{1/p_n}, \\
  \mathrm{S}_n^{\mathrm{inv}}(A,p_0) &:=
   \Big(\frac{C_{\mathrm{MA}}(d)}{(r_n-r_{n+1})^{2}}\,
     \Big(\Lambda\Big(\frac{p_n}{1+p_n}\Big)^{2}+1\Big)\Big)^{1/p_n}.
\end{align*}
The constant
\begin{equation*}
  C_{\mathrm{wH}}^{0}(d) :=
  \begin{cases}
    C_{\mathrm{Moser}}(d)\,(\chi^{2})^{\sum_{n\geq 0} n q^{n}}, & d > 2,\\
    C_{\mathrm{Moser}}(d), & \text{otherwise},
  \end{cases}
\end{equation*}
(\leanname{C\_weakHarnack0}) absorbs the extra geometric factor coming from the
$(\chi^{2})^{i}$ term in the inverse step bound.
\end{definition}

\begin{lemma}[Step bound]
\label{lem:stepInv-le}
\lean{DeGiorgi.superStepConstInv_le}
\leanok
\uses{def:ellipticityRatio, def:superIterIntegralInv}
For $d>2$ and $p_0 > 0$,
\begin{equation*}
  \mathrm{S}_n^{\mathrm{inv}}(A,p_0)
  \;\leq\;
  \Big(16\,C_{\mathrm{MA}}(d)\,(\Lambda p_0^{2}+1)\,(4\chi^{2})^{n}\Big)^{1/p_n}.
\end{equation*}
\end{lemma}

\begin{proof}
\leanok
\uses{def:C_MoserAnchor, lem:one_le_C_MoserAnchor, def:moserChi, lem:moserDecayRatio_nonneg, def:moserRadius, lem:moserRadius_zero, def:moserExponentSeq, lem:moserExponentSeq_zero}

The factor $C_{\mathrm{MA}}/(r_n - r_{n+1})^{2} = 16\,C_{\mathrm{MA}}\,4^{n}$ as
before. Using $p_n/(1+p_n) \leq p_n$ and $p_n^{2} = p_0^{2}\chi^{2n}$,
\begin{equation*}
  \Lambda\Big(\frac{p_n}{1+p_n}\Big)^{2}+1
  \;\leq\; \Lambda p_n^{2}+1
  \;\leq\; (\Lambda p_0^{2}+1)\,(\chi^{2})^{n}.
\end{equation*}
Multiplying gives the inner expression $16 C_{\mathrm{MA}}(\Lambda p_0^{2}+1)
(4\chi^{2})^{n}$; raising to $1/p_n$ yields the claim.
\end{proof}

\begin{theorem}[Inverse iteration]
\label{thm:iter-inv}
\lean{DeGiorgi.supersolution_iteration_inverse}
\leanok
\uses{def:ellipticityRatio, def:IsSubsolution, def:moserRadius, def:moserExponentSeq, def:superIterIntegralInv}
For $d>2$, $u>0$ a supersolution with $\int_{\Bz{1}} |u^{-1}|^{p_0}\,dx < \infty$,
and every $n \in \N$, $|u^{-1}|^{p_n}$ is integrable on $\Bz{r_n}$ and
\begin{equation*}
  N_n^{\mathrm{inv}}(u)
  \;\leq\;
  \prod_{i=0}^{n-1}\mathrm{S}_i^{\mathrm{inv}}(A,p_0)\;\cdot\; N_0^{\mathrm{inv}}(u).
\end{equation*}
\end{theorem}

\begin{proof}
\leanok
\uses{def:C_MoserAnchor, lem:one_le_C_MoserAnchor, def:moserChi, lem:moserRadius_zero, lem:moserExponentSeq_zero, thm:preMoser-inv}

Induction on $n$. The base case $n = 0$ uses $r_0 = 1$
(\leanname{moserRadius\_zero}) and $p_0 = p_0$
(\leanname{moserExponentSeq\_zero}): the integrability hypothesis for
$|u^{-1}|^{p_0}$ on $\Bz{1}$ is exactly what is assumed, and the empty product
is $1$ so the iteration bound reads $N_0^{\mathrm{inv}}(u) \le N_0^{\mathrm{inv}}(u)$.

For the inductive step from $n$ to $n+1$, set $p_n := \mathrm{moserExponentSeq}\,d\,p_0\,n$
and observe $p_n > 0$ by \leanname{moserExponentSeq\_pos} (using $d > 2$ and
$p_0 > 0$). Apply the inverse-power one-step pre-estimate
Theorem~\ref{thm:preMoser-inv}
(\leanname{supersolution\_preMoser\_inverse}) with parameters $(p, r, s) =
(p_n, r_{n+1}, r_n)$. The hypotheses required are: positivity of $p_n$ (from
the previous remark), positivity of $r_{n+1}$ (\leanname{moserRadius\_pos}),
the strict inequality $r_{n+1} < r_n$ (\leanname{moserRadius\_succ\_lt}), the
upper bound $r_n \le 1$ (\leanname{moserRadius\_le\_one}), the strict
positivity of $u$ on the unit ball, the supersolution property of $u$, and
the integrability of $|u^{-1}|^{p_n}$ on $\Bz{r_n}$ (which is the inductive
hypothesis after rewriting via $\mathrm{moserExponentSeq}\,d\,p_0\,n = p_n$).

The pre-estimate produces:
(i) integrability of $|u^{-1}|^{\chi p_n}$ on $\Bz{r_{n+1}}$, which by
$\chi p_n = p_{n+1}$ (\leanname{moserExponentSeq\_succ}) is exactly the
integrability assertion at level $n+1$; and
(ii) the one-step bound
$N_{n+1}^{\mathrm{inv}}(u) \le \mathrm{S}_n^{\mathrm{inv}}(A, p_0)\,
N_n^{\mathrm{inv}}(u)$.
Multiplying the inductive bound for $N_n^{\mathrm{inv}}$ by
$\mathrm{S}_n^{\mathrm{inv}}$ and absorbing the additional factor into the
range product
$\prod_{i=0}^{n} \mathrm{S}_i^{\mathrm{inv}} = \bigl(\prod_{i=0}^{n-1} \mathrm{S}_i^{\mathrm{inv}}\bigr)
\cdot \mathrm{S}_n^{\mathrm{inv}}$ closes the induction.
\end{proof}

\begin{lemma}[Geometric majorant]
\label{lem:geom-maj-inv}
\lean{DeGiorgi.supersolution_geometric_majorant_inv}
\leanok
\uses{def:ellipticityRatio, def:superIterIntegralInv}
For $d>2$ and $p_0>0$,
\begin{equation*}
  \prod_{i=0}^{n-1}\mathrm{S}_i^{\mathrm{inv}}(A,p_0)
  \;\leq\;
  \big(C_{\mathrm{wH}}^{0}(d)\,(\Lambda p_0^{2}+1)^{d/2}\big)^{1/p_0}.
\end{equation*}
\end{lemma}

\begin{proof}
\leanok
\uses{def:C_MoserAnchor, lem:one_le_C_MoserAnchor, def:moserChi, lem:moserDecayRatio_nonneg, def:C_Moser, def:moserExponentSeq, lem:moserExponentSeq_zero, lem:stepInv-le}

By Lemma~\ref{lem:stepInv-le}, the product is bounded by
\begin{equation*}
  \prod_{i=0}^{n-1}\big(K\,L\,(4\chi^{2})^{i}\big)^{1/p_i},
\end{equation*}
where $K := 16\,C_{\mathrm{MA}}(d)$, $L := \Lambda p_0^{2}+1$, $1/p_i = q^{i}/p_0$.
Taking logs and using $\sum_{i<n} q^{i} \leq d/2$ (\leanname{moserDecayRatio\_sum\_le\_half\_dim})
and $\sum_{i<n} i q^{i} \leq \sum_{i\geq 0} i q^{i}$
(\leanname{moserDecayRatio\_weighted\_sum\_le\_tsum}), the product is at most
\begin{equation*}
  \big(K^{d/2}\,L^{d/2}\,(4\chi^{2})^{\sum_{i\geq 0} i q^{i}}\big)^{1/p_0}.
\end{equation*}
Since $K^{d/2}\,4^{\sum_{i\geq 0} i q^{i}} \leq C_{\mathrm{Moser}}(d)$ by definition
of $C_{\mathrm{Moser}}$, and $(\chi^{2})^{\sum_{i\geq 0} i q^{i}}$ is exactly the
extra factor included in $C_{\mathrm{wH}}^{0}(d)$, the product is bounded by
$\big(C_{\mathrm{wH}}^{0}(d)\,L^{d/2}\big)^{1/p_0}$.
\end{proof}

\begin{theorem}[Inverse a.e.\ closeout]
\label{thm:ae-closeout-inv}
\lean{DeGiorgi.supersolution_ae_closeout_inv}
\leanok
\uses{def:ellipticityRatio, def:IsSubsolution, def:moserChi, def:C_Moser}
For $d>2$, $p_0 > 0$, $u>0$ a supersolution and $|u^{-1}|^{p_0}$ integrable on
$\Bz{1}$, for almost every $x \in \Bz{1/2}$,
\begin{equation*}
  |u(x)^{-1}|^{p_0}
  \;\leq\;
  C_{\mathrm{wH}}^{0}(d)\,(\Lambda p_0^{2}+1)^{d/2}\,
  \int_{\Bz{1}} |u^{-1}|^{p_0}\,dx.
\end{equation*}
\end{theorem}

\begin{proof}
\leanok
\uses{lem:moserDecayRatio_nonneg, lem:C_MoserAnchor_le_C_Moser, def:moserRadius, lem:moserRadius_zero, def:moserExponentSeq, lem:moserExponentSeq_zero, lem:moser-closeout, def:superIterIntegralInv, thm:iter-inv, lem:geom-maj-inv}

Combine Theorem~\ref{thm:iter-inv} with Lemma~\ref{lem:geom-maj-inv} to get,
for every $n$,
\begin{equation*}
  N_n^{\mathrm{inv}}(u)
  \;\leq\;
  \big(C_{\mathrm{wH}}^{0}(d)\,(\Lambda p_0^{2}+1)^{d/2}\,I_0^{\mathrm{inv}}(u)\big)^{1/p_0}.
\end{equation*}
The Markov-type closeout argument of Lemma~\ref{lem:moser-closeout}, applied
verbatim with $u_+$ replaced by $u^{-1}$ (and $1<p_0$ replaced by $p_0>0$, which
only affects the choice of test exponent $q_n = 2\chi^{n}$ but not the
convergence) yields the desired a.e.\ inequality.
\end{proof}

\begin{theorem}[Stage one weak Harnack (inverse)]
\label{thm:wH-stage1-inv}
\lean{DeGiorgi.weak_harnack_stage_one_inverse}
\leanok
\uses{def:ellipticityRatio, def:IsSubsolution, def:moserChi, def:C_Moser}
Let $d>2$, $A$ on $\Bz{1}$, $u>0$ on $\Bz{1}$ a supersolution, and $p_0 > 0$. Then
\begin{multline*}
  (\Lambda p_0^{2}+1)^{-d/2}\,
  \Big(\int_{\Bz{1}} |u^{-1}|^{p_0}\,dx\Big)^{-1}
  \\
  \;\leq\; C_{\mathrm{wH}}^{0}(d)\;(\essinf_{\Bz{1/2}} u)^{p_0}.
\end{multline*}
\end{theorem}

\begin{proof}
\leanok
\uses{lem:moserDecayRatio_nonneg, lem:C_MoserAnchor_le_C_Moser, thm:ae-closeout-inv, lem:meas-bd}

If $|u^{-1}|^{p_0}$ is not integrable on $\Bz{1}$, the integral is taken to be
$+\infty$ and the LHS is $0 \leq \mathrm{RHS}$.

Otherwise, set $I := \int_{\Bz{1}} |u^{-1}|^{p_0}\,dx$ and
$c_0 := C_{\mathrm{wH}}^{0}(d)\,(\Lambda p_0^{2}+1)^{d/2}\,I$. By
Theorem~\ref{thm:ae-closeout-inv} we have $|u^{-1}|^{p_0} \leq c_0$ a.e.\ on
$\Bz{1/2}$, equivalently $u^{p_0} \geq c_0^{-1}$ a.e., i.e.\
$u \geq c_0^{-1/p_0}$ a.e.\ on $\Bz{1/2}$. Therefore
$\essinf_{\Bz{1/2}} u \geq c_0^{-1/p_0}$, so
$(\essinf_{\Bz{1/2}} u)^{p_0} \geq c_0^{-1}$. Multiplying by
$C_{\mathrm{wH}}^{0}(d) \geq 0$ and using
$L^{-d/2} I^{-1} = C_{\mathrm{wH}}^{0}(d)\,c_0^{-1}$ (where $L = \Lambda p_0^{2}+1$)
gives the displayed inequality.
\end{proof}

\subsection{Stage two (forward low-power iteration)}

The forward iteration is run only finitely many steps: starting from
$p_0 := q\,\chi^{-m}$ for a target exponent $q \in (0,1)$ and a depth $m$
chosen so that $p_0 \leq p$, we iterate $m+1$ times to reach the level $p_m =
q$. The forward step constant
\begin{equation*}
  \mathrm{S}_n^{\mathrm{fwd}}(A,p_0) :=
  \Big(\frac{C_{\mathrm{MA}}(d)}{(r_n-r_{n+1})^{2}}\,
    \Big(\Lambda\Big(\frac{p_n}{1-p_n}\Big)^{2}+1\Big)\Big)^{1/p_n}
\end{equation*}
\leanname{superStepConstFwd} is finite as long as $p_n < 1$, which is
guaranteed by $p_n \leq q < 1$ for $n \leq m$.

\begin{lemma}[Forward depth]
\label{lem:fwd-depth}
\lean{DeGiorgi.exists_forward_iteration_depth}
\leanok
For $0 < p < q < 1$ there exists $m \in \N$ with
\begin{equation*}
  q\,\chi^{-m} < p \quad \text{and} \quad p \leq \chi\,\big(q\,\chi^{-m}\big).
\end{equation*}
\end{lemma}

\begin{proof}
\leanok
\uses{def:moserChi, lem:moserDecayRatio_nonneg}

Since $\chi > 1$, $q\,\chi^{-m} \to 0$ as $m \to \infty$, so any sufficiently
large $m$ satisfies the first inequality; choose the smallest such $m$, then
$q\,\chi^{-(m-1)} \geq p$, equivalently $p \leq \chi\,(q\,\chi^{-m})$.
\end{proof}

\begin{lemma}[Forward step bound]
\label{lem:superStepConstFwd_le}
\lean{DeGiorgi.superStepConstFwd_le}
\leanok
\uses{def:ellipticityRatio}
For $d>2$, $0 < q < 1$, $m \in \N$, $i \leq m$, and $p_0 := q\,\chi^{-m}$,
\begin{equation*}
  \mathrm{S}_i^{\mathrm{fwd}}(A,p_0)
  \;\leq\;
  \Big(16\,C_{\mathrm{MA}}(d)\,\Big(\frac{\Lambda p_0^{2}}{(1-q)^{2}}+1\Big)\,
  (4\chi^{2})^{i}\Big)^{1/p_i}.
\end{equation*}
\end{lemma}

\begin{proof}
\leanok
\uses{def:C_MoserAnchor, lem:one_le_C_MoserAnchor, def:moserChi, lem:moserDecayRatio_nonneg, def:moserRadius, lem:moserRadius_zero, def:moserExponentSeq, lem:moserExponentSeq_zero}

For $i \leq m$, $p_i = p_0\,\chi^{i} = q\,\chi^{i-m} \leq q$, hence
$1 - p_i \geq 1 - q > 0$ and
\begin{equation*}
  \frac{p_i}{1-p_i} \leq \frac{p_i}{1-q}, \qquad
  \Big(\frac{p_i}{1-p_i}\Big)^{2}\leq \frac{p_i^{2}}{(1-q)^{2}}
  \leq \frac{p_0^{2}\chi^{2i}}{(1-q)^{2}}.
\end{equation*}
Inserting into $\mathrm{S}_i^{\mathrm{fwd}}$ and using
$C_{\mathrm{MA}}/(r_i - r_{i+1})^{2} = 16\,C_{\mathrm{MA}}\,4^{i}$ gives the bound.
\end{proof}

\begin{theorem}[Forward iteration]
\label{thm:iter-fwd}
\lean{DeGiorgi.supersolution_iteration_forward}
\leanok
\uses{def:ellipticityRatio, def:IsSubsolution}
For $d>2$, $u>0$ a supersolution, $0<q<1$, $m \in \N$, set $p_0 := q\,\chi^{-m}$.
If $|u|^{p_0}$ is integrable on $\Bz{1}$, then for every $n \leq m+1$
$|u|^{p_n}$ is integrable on $\Bz{r_n}$ and
\begin{equation*}
  N_n^{\mathrm{fwd}}(u)
  \;\leq\;
  \prod_{i=0}^{n-1} \mathrm{S}_i^{\mathrm{fwd}}(A,p_0)\;\cdot\;
  N_0^{\mathrm{fwd}}(u).
\end{equation*}
\end{theorem}

\begin{proof}
\leanok
\uses{def:C_MoserAnchor, lem:one_le_C_MoserAnchor, def:moserChi, lem:moserDecayRatio_nonneg, def:moserRadius, lem:moserRadius_zero, def:moserExponentSeq, lem:moserExponentSeq_zero, thm:preMoser-fwd}

Induction on $n \leq m+1$. The base case is trivial. For the step, note that
$p_n = p_0\,\chi^{n} = q\,\chi^{n-m}$, so for $n \leq m$ we have $p_n \leq q < 1$
(this is what allows the application of Theorem~\ref{thm:preMoser-fwd}). The
one-step bound gives $N_{n+1}^{\mathrm{fwd}}(u) \leq
\mathrm{S}_n^{\mathrm{fwd}}(A,p_0)\,N_n^{\mathrm{fwd}}(u)$, and combining with
the inductive hypothesis yields the displayed product bound.
\end{proof}

\begin{lemma}[Forward geometric majorant]
\label{lem:geom-maj-fwd}
\lean{DeGiorgi.supersolution_geometric_majorant_fwd}
\leanok
\uses{def:ellipticityRatio}
With the notation above, for any $n \leq m+1$,
\begin{equation*}
  \prod_{i=0}^{n-1}\mathrm{S}_i^{\mathrm{fwd}}(A,p_0)
  \;\leq\;
  \Big(C_{\mathrm{wH}}^{0}(d)\,
    \Big(\frac{\Lambda p_0^{2}}{(1-q)^{2}}+1\Big)^{d/2}\Big)^{1/p_0}.
\end{equation*}
\end{lemma}

\begin{proof}
\leanok
\uses{def:C_MoserAnchor, lem:one_le_C_MoserAnchor, def:moserChi, lem:moserDecayRatio_nonneg, def:C_Moser, def:moserRadius, def:moserExponentSeq, lem:moserExponentSeq_zero, lem:geom-maj-inv, lem:superStepConstFwd_le}

Identical to the proof of Lemma~\ref{lem:geom-maj-inv}, with $\Lambda p_0^{2}+1$
replaced by $\Lambda p_0^{2}/(1-q)^{2}+1$ throughout.
\end{proof}

\begin{definition}[Forward closeout constant]
\label{def:C_weakHarnack0Forward}
\lean{DeGiorgi.C_weakHarnack0Forward}
\leanok
\uses{def:moserChi, def:C_Moser}
For $d>2$,
\begin{equation*}
  C_{\mathrm{wH,fwd}}^{0}(d) :=
  \big(C_{\mathrm{wH}}^{0}(d)\,(|\Bz{1}|+1)\big)^{\chi}.
\end{equation*}
\end{definition}

\begin{theorem}[Stage one weak Harnack (forward)]
\label{thm:wH-stage1-fwd}
\lean{DeGiorgi.weak_harnack_stage_one_forward}
\leanok
\uses{def:ellipticityRatio, def:IsSubsolution, def:C_weakHarnack0Forward}
Let $d>2$, $A$ on $\Bz{1}$, $u>0$ a supersolution on $\Bz{1}$, and
$0 < p < q < 1$. Then
\begin{multline*}
  \Big(\int_{\Bz{1/2}} |u|^{q d/(d-2)}\,dx\Big)^{p(d-2)/(qd)}
  \\
  \;\leq\;
  C_{\mathrm{wH,fwd}}^{0}(d)\,
  \Big(\frac{\Lambda p^{2}}{(1-q)^{2}}+1\Big)^{d\chi/2}
  \int_{\Bz{1}} |u|^{p}\,dx.
\end{multline*}
\end{theorem}

\begin{proof}
\leanok
\uses{def:MemW1p, def:MemW1pWitness, def:C_MoserAnchor, def:moserChi, lem:moserDecayRatio_nonneg, def:C_Moser, lem:C_MoserAnchor_le_C_Moser, def:moserRadius, lem:moserRadius_zero, def:moserExponentSeq, lem:moserExponentSeq_zero, lem:fwd-depth, thm:iter-fwd, lem:geom-maj-fwd}

Set $\chi := d/(d-2)$, $q\chi := q d/(d-2)$. By Lemma~\ref{lem:fwd-depth} pick
$m$ with $p_0 := q\chi^{-m} < p$ and $p \leq \chi p_0$, and note $p_0 < p < 1$.
Theorem~\ref{thm:iter-fwd} with this $p_0$ at depth $n = m+1$ gives the iterated
$L^{p_{m+1}} = L^{q\chi}$ bound on $\Bz{r_{m+1}}$, namely
\begin{multline*}
  \Big(\int_{\Bz{r_{m+1}}} |u|^{q\chi}\,dx\Big)^{1/(q\chi)}
  \\ \leq\;
  \prod_{i=0}^{m}\mathrm{S}_i^{\mathrm{fwd}}(A,p_0)\;\cdot\;
  \Big(\int_{\Bz{1}} |u|^{p_0}\,dx\Big)^{1/p_0}.
\end{multline*}
By Lemma~\ref{lem:geom-maj-fwd}, the product is at most
\begin{equation*}
  \Big(C_{\mathrm{wH}}^{0}(d)\,
    \big(\Lambda p_0^{2}/(1-q)^{2}+1\big)^{d/2}\Big)^{1/p_0}.
\end{equation*}
Use $\Bz{1/2} \subseteq \Bz{r_{m+1}}$ to restrict to $\Bz{1/2}$, raise the
inequality to the $p$-th power, and absorb the change of base from $p_0$ to $p$
via the volume comparison
\leanname{weak\_harnack\_stage\_one\_forward\_power\_upgrade}: by Holder with
exponents $p_0 < p \leq \chi p_0$ on the unit ball,
\begin{equation*}
  \Big(\int_{\Bz{1}} |u|^{p_0}\,dx\Big)^{1/p_0}
  \leq
  \Big(\int_{\Bz{1}} |u|^{p}\,dx\Big)^{1/p}\,
  |\Bz{1}|^{1/p_0 - 1/p},
\end{equation*}
and consequently
\begin{equation*}
  \Big(\int_{\Bz{1}} |u|^{p_0}\,dx\Big)^{p/p_0}
  \leq (|\Bz{1}|+1)^{\chi}\,\int_{\Bz{1}} |u|^{p}\,dx,
\end{equation*}
which absorbs the volume into $C_{\mathrm{wH,fwd}}^{0}(d)$. Together with
$(\Lambda p_0^{2}/(1-q)^{2}+1)^{d/(2)\cdot p/p_0} \leq
(\Lambda p^{2}/(1-q)^{2}+1)^{d\chi/2}$
(monotonicity in both base and exponent, since $p/p_0 \leq \chi$ and
$p_0^{2} \leq p^{2}$), the displayed inequality follows. The exponent
$p(d-2)/(qd) = p/(q\chi)$ comes from raising the $L^{q\chi}$-norm to the
$p$-th power.
\end{proof}

\begin{remark}[Reverse-Holder reformulation]
Theorems~\ref{thm:wH-stage1-inv} and~\ref{thm:wH-stage1-fwd} are precisely the
two halves of the weak Harnack inequality at this stage of the development:
the first asserts a lower bound on $u$ in $\Bz{1/2}$ in terms of an inverse
integral on $\Bz{1}$, while the second is a forward $L^{q\chi}$--$L^{p}$
reverse Holder inequality on the same pair of balls. The crossover between
the two stages, completing the weak Harnack inequality, is performed in the
file \texttt{WeakHarnack.lean} using the constants $C_{\mathrm{wH}}^{0}(d)$
and $C_{\mathrm{wH,fwd}}^{0}(d)$ assembled here.
\end{remark}


%% ========================================================================
